\begin{table}[ht!]
\centering
\caption{The level after adoption, read against the months before the rate rise.}
\label{tab:window}
\footnotesize
\begin{tabular}{@{}lcc@{}}
\toprule
Term & 22--25 & 26--30 \\
\midrule
Tightening months, April to November 2022 & $+0.0159^{*}$ (0.0075) & $+0.0184^{*}$ (0.0044) \\
The thirteen months after the launch & $-0.0020$ (0.0135) & $+0.0104$ (0.0080) \\
The level after adoption, from January 2024 & $-0.0419^{*}$ (0.0189) & $-0.0298^{*}$ (0.0123) \\
The step from the 2023 level & $-0.0399^{*}$ (0.0102) & $-0.0403^{*}$ (0.0067) \\
Employers & 104{,}217 & 117{,}090 \\
\bottomrule
\end{tabular}
\begin{minipage}{0.88\textwidth}\footnotesize\vspace{4pt}
Poisson pseudo-maximum likelihood on employer-by-age-by-month counts, with employer-by-month, employer-by-age and month-by-age effects; exposure is the employer's 2019 occupation mix; standard errors clustered by employer in parentheses. The calendar-quarter terms are in the fit and are not printed. Here the Riksbank interaction is a window, April to November 2022, rather than an indicator left switched on through the post period, so the second and third rows are levels against January 2021 to March 2022 and not steps. The two specifications describe the same fitted means, the window estimate being the sum of the cumulative ones exactly. The fourth row is the third minus the second, with the standard error from the covariance of the two terms; it is the step Table~1 of the paper reports. Each band is fitted on its own panel. $^{*}$ $p<0.05$. Source: script 83.
\end{minipage}
\end{table}
